\section{Part A: disposition and scope}
\label{A:unresolved}
Every requested Part A assertion is now decided in its stated
scope. The counterexamples and their proved replacements remain
separate from further classification questions about special fields.

\subsection{A.2 is refuted with persistent failure on every open}
Theorem~\ref{A:A2-persistent-arithmetic-counterexample} constructs
an excellent regular Dedekind domain $R$ with $2$ invertible and
$\cd_2(\Frac R)=1$, but $\cd_2(U)=2$ for every nonempty affine
open $U$. Its Picard group modulo two is infinite dimensional
after every finite deletion. The quadratic-step construction
preserves old classes modulo two by an explicit diagonal transfer,
while controlling splitting and ramification at every old prime.
Thus both the quasi-excellent and unrestricted readings of A.2
are refuted. Corollary~\ref{A:A2-sharp-curve-consequences} also
gives persistent dimension-one refutations of A.1 and both A.6
claims, with no orderable residue fields.

The general lower witness after shrinking remains valid by
Proposition~\ref{A:spread}. On curves, the exact replacement is
the simultaneous Milnor-residue criterion of
Theorems~\ref{A:Dedekind-Milnor-criterion} and
\ref{A:curve-normalization}, with the sharp upper bound one
greater than the generic dimension. This covers nonnormal and
nonexcellent curves. In generic dimension one the criterion is
divisibility modulo $p$ of the Picard groups of all finite
separable normalizations after one localization.

The positive results proved during the investigation retain their
full statements: all semilocal curves; schemes essentially of
finite type over a field; every noetherian model of an arithmetic
field of maximal cohomological dimension
$\operatorname{trdeg}_{\Q}K+2$; all models of the specified
strictly local surface fields; and semilocal or topologically
normal models of quasi-excellent strictly local fraction fields.
The Nagata--Gabber example satisfies eventual equality after its
specified parameter localization.

The cyclotomic comparison in
Theorems~\ref{A:Iwasawa-equivalence} and
\ref{A:canonical-Iwasawa} is still an exact equivalence for that
particular family. The non-Galois counterexample neither proves
nor refutes the full cyclotomic $\mu$-conjecture. Settling each
cyclotomic instance is not required to refute the universal A.2
assertion, and is not counted as remaining requested work.

\subsection{The other Part A questions}
Theorem~\ref{A:wild-affine-counterexample} supplies the excellent
regular complete two-dimensional counterexample to A.5, and
Theorem~\ref{A:wild-affine-upper} proves its uniform replacement.
The affine A.7 assertions hold over quasi-excellent bases and
fail without that hypothesis; their nonaffine versions as posed
fail without separatedness. The fraction-field and local lower
bounds in A.8 and A.9 hold in their precisely stated scopes,
with the real and characteristic exceptions explicitly decided.
The outcome table lists all individual variants and proof locators.

\begin{remark}[External modern reading consulted]
We consulted Michel Hua's project, including the November 2026
modernized-reading appendix \citep[pp.~42--43]{A_Hua2026} and its
separate conjectures article. Its attention to the distinction between
the open $D(p)$ and a general mixed-characteristic affine scheme is
consistent with the scope checks above. We credit that discussion to
Michel Hua's project and have independently checked the cited inputs.
One source attribution in the conjectures article needs narrowing:
\citet[Remark 7.4]{A_BhattMathew2020} discusses topologically finite-type
algebras over an algebraically closed complete nonarchimedean field.
It does not state that the arbitrary global mixed-characteristic
scheme bound of 31/085 is open. We use neither that remark nor Hua's
status description as evidence for a theorem or for its failure.
Our searches located this modern project but did not establish the
absence of other studies; no bibliographical exhaustiveness or
priority claim is made.
\end{remark}

